\documentclass[11pt,a4paper]{article}

\usepackage[utf8]{inputenc}
\usepackage[T1]{fontenc}
\usepackage{amsmath,amssymb}
\usepackage{graphicx}
\usepackage{booktabs}
\usepackage{hyperref}
\usepackage{xcolor}
\usepackage{tcolorbox}
\usepackage[margin=25mm]{geometry}
\usepackage{xeCJK}
\setCJKmainfont{Hiragino Mincho ProN}
\setCJKsansfont{Hiragino Sans}

\title{商用 SMR/FBAR による DN 真空効果の検証\\
\large ポケットの中の Boyer 予言を暴く}
\author{Wright Brothers Research Program}
\date{2026年4月}

\begin{document}
\maketitle

\begin{abstract}
本ガイドは Wright Brothers プロジェクトにおける DN カシミール検証の
\textbf{最速ルート}として，商用 RF フィルタ産業が大量生産している
SMR (Solidly Mounted Resonator) と FBAR (Film Bulk Acoustic
Resonator) の比較実験を詳述する．SMR は構造上 Dirichlet--Neumann
境界を真に実装しており，FBAR は Neumann--Neumann 境界．両者は
AlN 薄膜という同一材料で作られ，世界中の 5G 携帯電話に毎日
数十億個搭載されている．つまり\textbf{Boyer (1974) が予言した
DN 真空構造は，既に人類のポケットの中に存在している}．本ガイドでは，
物理原理，入手方法（商用評価キット $\sim$\$500），測定プロトコル，
期待される信号強度，そして論文化までの最短経路を提示する．総額
30--100 万円で週末から着手可能．cQED ルートや BAW ルートより
数桁安く，かつ $\zeta_{\neg 2}$ 予測の核心を直接ついている．
\end{abstract}

\tableofcontents

\newpage
\section{なぜ SMR/FBAR が最速ルートなのか}

\subsection{これまでの選択肢の整理}

Wright Brothers プロジェクトが検討してきた DN カシミール検証ルート：

\begin{center}
\begin{tabular}{llll}
\toprule
ルート & 境界条件 & 予算 & 状況 \\
\midrule
軸 A: cQED ($\lambda/4$ vs $\lambda/2$) & 電磁 DN & 数千万円 & 施設必須 \\
軸 A: BAW (水晶, クランプ貼付) & 機械 DN & 100 万円 & 自作可能 \\
軸 B: 音響アナログ & フォノン DN & 5 万円 & 古典領域のみ \\
\textbf{本ガイド: SMR/FBAR 比較} & 機械 DN & 30 万円 & 商用品利用 \\
\bottomrule
\end{tabular}
\end{center}

SMR/FBAR ルートの特徴：
\begin{itemize}
\item \textbf{商用品ベース}：自分で DN 境界を作る必要がない
\item \textbf{量産品質}：半導体プロセスで作られており，製造バラつきが
極小
\item \textbf{既存評価インフラ}：RF/MEMS 研究室には必ず測定系がある
\item \textbf{産業界の既存データ}：S パラメータ測定が膨大に蓄積
（ただし物理学的視点での再解釈は未実施）
\end{itemize}

\subsection{一番大事な事実}

5G 携帯電話のバンドフィルタに使われている SMR は構造上
\textbf{Dirichlet--Neumann 境界を実装しており}，これは Boyer (1974) が
50 年前に予言した真空エネルギー符号反転の構造と数学的に同型である．

\begin{tcolorbox}[colback=yellow!10,colframe=orange!60!black,title=もし我々の仮説が正しいなら]
SMR 内部の AlN 薄膜は，FBAR 内部の同じ AlN 薄膜と比べて，
\textbf{真空エネルギーの符号が異なる}はずである．この差は：
\begin{itemize}
\item 基本共振周波数のわずかなシフト
\item 非線形アンハーモニシティの符号差
\item $Q$ 値の周波数依存性の違い
\end{itemize}
として現れる．\textbf{これは誰も測定していない}．産業界は
S$_{11}$ と挿入損失しか測らないから．
\end{tcolorbox}

\section{物理：SMR と FBAR の構造と境界条件}

\subsection{FBAR の構造}

Film Bulk Acoustic Resonator：

\begin{center}
\begin{tabular}{|c|}
\hline
上部電極 (Mo or W, 100--300\,nm) \\
\hline
\textbf{AlN 圧電膜 ($d \approx 1\,\mu$m)} \\
\hline
下部電極 (Mo or W, 100--300\,nm) \\
\hline
空気（エアギャップ 2--5\,$\mu$m）\\
\hline
Si 基板 \\
\hline
\end{tabular}
\end{center}

下部電極の下にエアギャップがあり，AlN は基板から「宙吊り」状態．
両面とも実質的に空気/真空に接する．

\textbf{境界条件}：
\begin{itemize}
\item 上面：薄い金属電極＋空気．電極の質量負荷は微小なので
$\partial u/\partial z \approx 0$，\textbf{Neumann}
\item 下面：薄い金属電極＋空気（ギャップ）．同様に \textbf{Neumann}
\end{itemize}

NN 境界 → モード: $u_n(z) = \cos(n\pi z/d)$, $n=0,1,2,\ldots$\\
周波数: $f_n = n v_s/(2d)$，\textbf{全倍音（$n=1,2,3,4,\ldots$）}

実際には $n=0$ モード（一様変位）は物理的に意味がないので $n=1$
から始まる．

\subsection{SMR の構造}

Solidly Mounted Resonator：

\begin{center}
\begin{tabular}{|c|}
\hline
上部電極 (Mo or W) \\
\hline
\textbf{AlN 圧電膜 ($d \approx 1\,\mu$m)} \\
\hline
下部電極 (Mo or W) \\
\hline
\textbf{音響 Bragg 反射鏡（W/SiO$_2$ 多層膜，3--5 ペア）}\\
\hline
Si 基板 \\
\hline
\end{tabular}
\end{center}

AlN 圧電膜の下に\textbf{音響 Bragg 反射鏡}（DBR と同じ概念，ただし
音響波用）が配置されている．Bragg 反射鏡は W（高インピーダンス
$Z_W \approx 100$ MRayl）と SiO$_2$（低インピーダンス $Z_{\rm SiO_2}
\approx 13$ MRayl）の $\lambda/4$ 層を交互に積層する．

インピーダンスコントラスト $Z_W/Z_{\rm SiO_2} \approx 7.7$ は非常に
大きく，3 ペア で反射率 $R > 0.99$，5 ペアで $R > 0.9999$ に達する．
これは\textbf{ほぼ完全な Dirichlet 境界}．

\textbf{境界条件}：
\begin{itemize}
\item 上面：空気 → \textbf{Neumann}（FBAR と同じ）
\item 下面：Bragg 反射鏡 → \textbf{Dirichlet}（反射率 $>$ 0.999）
\end{itemize}

DN 境界 → モード: $u_n(z) = \cos((2n-1)\pi z/(2d))$, $n=1,2,\ldots$\\
周波数: $f_n = (2n-1) v_s/(4d)$，\textbf{奇数倍音のみ}

\subsection{構造の比較}

\begin{center}
\begin{tabular}{lll}
\toprule
 & FBAR & SMR \\
\midrule
上面境界 & 自由 (N) & 自由 (N) \\
下面境界 & 自由 (N) & Bragg (D) \\
モード & 全倍音 $n=1,2,3,\ldots$ & 奇数 $2n-1=1,3,5,\ldots$ \\
基本周波数 $f_0$ & $v/(2d)$ & $v/(4d)$ \\
温度特性 & 中 & 良 \\
Q 値 (2.4 GHz, 室温) & $\sim 3000$ & $\sim 2500$ \\
スプリアスモード & 少 & 中 \\
製造歩留 & 低 & 高 \\
典型用途 & LTE/WiFi & 5G, 厳しい帯域 \\
\bottomrule
\end{tabular}
\end{center}

\textbf{同じ AlN 厚さなら SMR は FBAR の半分の基本周波数になる}ことに
注意．比較実験では「FBAR の奇数倍音位置 ($f_0, 3f_0, 5f_0,\ldots$)」と
「SMR の全モード ($(2n{-}1)f_0^{SMR}$)」を一致させる必要がある．

具体的には，FBAR の $d_{\rm FBAR}$ と SMR の $d_{\rm SMR} = d_{\rm FBAR}/2$
を用意すれば，両者の奇数倍音位置が完全に一致する．

\subsection{Bragg 反射鏡の品質を甘く見積もらないために}

SMR の Bragg 反射鏡は設計帯域内（通常 2--3 GHz 付近）でのみ
高反射率．この帯域外では反射率が落ち，DN 境界が崩れる．
具体的には：

\begin{itemize}
\item 設計中心周波数 $\pm 20\%$: $R > 0.999$（DN ほぼ完璧）
\item $\pm 50\%$: $R \sim 0.95$（DN 近似）
\item 2 倍以上離れる: $R < 0.5$（境界が不明瞭）
\end{itemize}

したがって DN 性質が保証されるのは設計中心周波数のある狭い帯域．
我々が見たい倍音 comb のどこまでを SMR の Bragg 鏡がカバーできるか，
事前にメーカーからデータシートを取って確認する必要がある．

\begin{tcolorbox}[colback=red!5,colframe=red!50!black,title=落とし穴]
「SMR は DN」という主張は\textbf{Bragg 鏡の帯域内でのみ}成立する．
帯域外では実質 NN に戻る．偶数倍音の「欠落」が設計帯域内のみで
起こり，帯域外では現れることを想定してデータ解析すべき．これ自体が
面白い検証点になる（Bragg 帯域内/外で効果のオン/オフが見える）．
\end{tcolorbox}

\section{先行研究と空白地帯}

\subsection{FBAR/SMR 研究の文脈}

FBAR/SMR は RF フィルタとして 2000 年代初頭に商用化（Aigner et al.\
2001, Agilent / Infineon）．以降，5G 帯域フィルタの主力．
物理的・工学的な研究は膨大：

\begin{itemize}
\item Ruby et al.\ (2002, 2007): Agilent/Avago の先駆的論文
\item Aigner et al.\ (2003): Infineon SMR の設計
\item Enzelberger et al.\ (2006): FBAR 温度特性
\item Hashimoto 編 \textit{RF Bulk Acoustic Wave Filters for Communications}
(Artech House, 2009): 標準教科書
\end{itemize}

これらは全て RF 工学視点．

\subsection{量子・真空物理学視点での空白}

\textbf{驚くべきことに}，SMR/FBAR を\textbf{真空エネルギーや
Casimir 効果の文脈で扱った論文は存在しない}（我々の調査範囲）．
理由：
\begin{itemize}
\item RF コミュニティは古典 S パラメータしか測らない
\item 量子音響コミュニティ（Chu, Cleland, Safavi-Naeini）は自前の
HBAR や optomechanical crystal を使い，産業品を使わない
\item Casimir コミュニティは電磁場に集中し，機械的類似物を見ない
\end{itemize}

この三分野が交わる場所に SMR/FBAR 真空比較実験がある．
\textbf{Wright Brothers が最初に踏み込む余地が十分に残されている}．

\subsection{関連する量子音響の先行研究}

DN ではないが，機械的真空を扱った重要な論文：

\begin{itemize}
\item O'Connell et al.\ (2010, Nature 464:697): 機械共振器の
基底状態冷却
\item Chu et al.\ (2017, 2018): HBAR 量子音響（これは NN 構造）
\item Arrangoiz-Arriola et al.\ (2019, Nature 571:537): フォノン
キャビティでの単一音子制御
\item Bienfait et al.\ (2019, Science 364:368): 量子音響伝送
\end{itemize}

\section{実験プロトコル}

\subsection{サンプル選定}

必要なペア：
\begin{itemize}
\item FBAR サンプル: AlN 厚さ $d_1$，基本周波数 $f_0 = v_{\rm AlN}/(2 d_1)$
\item SMR サンプル: AlN 厚さ $d_2 = d_1/2$，基本周波数 $f_0^{SMR} = v_{\rm AlN}/(4 d_2) = f_0$
\end{itemize}

$v_{\rm AlN} \approx 10400$\,m/s （縦波）なので，$f_0 = 2.6$\,GHz に
合わせるなら $d_1 = 2\,\mu$m, $d_2 = 1\,\mu$m．

\textbf{重要}: 同じメーカーの同じプロセスで作られたペアを選ぶ．
異なるメーカーだと AlN の $c$ 軸配向や応力が違い，$v_s$ が数\% ずれる．

\subsection{測定装置}

最低限必要：
\begin{itemize}
\item ベクトルネットワークアナライザ (VNA)，10\,MHz--8\,GHz：
中古 Keysight E5071B で ¥500,000，または nanoVNA V2 Plus4 で ¥25,000
\item RF プローブステーション（GSG プローブ）：中古 ¥200,000
\item 温度制御ステージ（ペルチェ，$-20$ から $+100^\circ$C）：¥50,000
\item ローノイズ RF 源（位相雑音測定用）：¥100,000
\item 測定 PC，校正キット，消耗品：¥150,000
\end{itemize}

合計 ¥1,000,000 程度．大学の RF 研究室なら全て揃っている．

\subsection{測定 1: モードスペクトル確認（基礎）}

\begin{enumerate}
\item FBAR と SMR を VNA に接続
\item $S_{11}$ (反射係数) を 100\,MHz--8\,GHz で掃引
\item 共振ディップの位置を抽出
\item 予測：
  \begin{itemize}
  \item FBAR: $f_0, 2f_0, 3f_0, 4f_0, 5f_0,\ldots$ のディップ
  \item SMR: $f_0, 3f_0, 5f_0, 7f_0,\ldots$ のディップ（偶数欠落）
  \end{itemize}
\item 偶数倍音の欠落を\textbf{有意な統計量}で確認（ノイズレベル対比
30 dB 以上を目標）
\end{enumerate}

これが\textbf{軸 D 測定 A と同型の基礎チェック}．ここで偶数倍音が
きれいに欠けていれば，DN 境界の実在が確定．

\subsection{測定 2: 精密周波数・Q 値測定}

\begin{enumerate}
\item 各共振ディップを中心に狭帯域スイープ（$\pm 10$\,MHz, 10000 点）
\item ローレンツ関数フィットで中心周波数 $f_n$ と Q 値を抽出
\item 複数回測定し Allan 分散を評価
\item 周波数精度 $10^{-6}$ を目標（2.6\,GHz なら $\pm 2.6$\,kHz）
\end{enumerate}

\subsection{測定 3: 真空 Lamb shift 探索（本丸）}

FBAR と SMR の基本周波数は，加工精度で決まる絶対値と，高次モード
との virtual coupling による Lamb shift の和：
\begin{equation}
  f_1^{\rm obs} = f_1^{(0)} + \sum_{n\geq 2} \frac{g_n^2}{f_1 - f_n}
\end{equation}
ここで $g_n$ はモード間の非線形結合．FBAR (NN) では全 $n$ の和，
SMR (DN) では奇数 $n$ のみの和．差：
\begin{equation}
  \Delta f = \sum_{n\;\text{even}} \frac{g_n^2}{f_1 - f_n}
\end{equation}

見積もり：$g/f_1 \sim 10^{-4}$, $N_{\rm even} \sim 3$ で
$\Delta f/f_1 \sim 10^{-8}$，$f_1 = 2.6$\,GHz なら $\Delta f \sim 26$\,Hz．

\textbf{問題}: FBAR と SMR は\textbf{別チップ}なので，加工誤差
$f_1^{(0)}$ の違いが kHz オーダーで存在し，Lamb shift 差 26\,Hz は
埋もれる．

解決策：\textbf{温度スイープ差分法}．

\subsection{測定 4: 温度スイープ差分法（系統誤差キラー）}

加工誤差は温度に依存しない．Lamb shift は高次モードとの間隔に
依存し，温度で変化する（モード全体が熱膨張でシフト）．

\begin{enumerate}
\item FBAR と SMR を同じ温度制御ステージに載せる
\item $-20^\circ$C から $+80^\circ$C まで $10^\circ$C 刻みで $f_1$ を測定
\item 温度係数 $df_1/dT$ を FBAR と SMR で比較
\item 純粋な材料の温度係数（両者共通）を差し引くと，モード構造
由来の温度依存性が残る
\item \textbf{FBAR と SMR の $df_1/dT$ の差}がDN/NN 差の現れ
\end{enumerate}

この差分法は加工誤差と共通系統を同時にキャンセルする．

\subsection{測定 5: 駆動振幅スイープ（アンハーモニシティ）}

強く駆動すると非線形効果で周波数がシフト（Duffing 振動子）：
\begin{equation}
  f_1(P) = f_1^{(0)} + \alpha P + \beta P^2 + \cdots
\end{equation}

アンハーモニシティ係数 $\alpha, \beta$ はモード結合 $g_n$ の関数．
FBAR と SMR で比較すると偶数モードの寄与が抽出できる．

駆動パワー $-20$\,dBm から $+10$\,dBm まで掃引．各パワーで $f_1$ と
$Q$ を記録．Duffing 理論でフィット．係数比較．

\subsection{測定 6: 低温測定（任意）}

液体窒素 77\,K で Q が一桁向上．信号の S/N が 10 倍になる．
温度制御ステージではなく液体窒素デュワーが必要（¥100,000 追加）．

\section{商用品の入手}

\subsection{FBAR を買う}

\subsubsection{評価キット}

\begin{itemize}
\item \textbf{Broadcom/Avago}: Acoustic Filter Evaluation Kit．
\$500--\$2,000．WiFi 6E / 5G 帯域．データシート付属．
\item \textbf{Qorvo}: TQ Series eval board．\$800 から．
\item \textbf{Murata}: SAW/BAW フィルタ評価キット．日本国内から
購入可能．¥80,000 から
\end{itemize}

\subsubsection{個別部品}

Digi-Key / Mouser で RF フィルタとして流通．1 個 \$5 から \$50．
デパッケージ（金属缶から剥き出しにする）が必要な場合は業者委託
（¥30,000 程度）．

\subsubsection{研究用サンプル}

学術目的で「未包装の bare die」を請求できる．問い合わせ先：
\begin{itemize}
\item Broadcom Innovation Research Lab (IRL) ：
研究用途で bare die 提供実績あり
\item Qorvo University Program ：年 2 回サンプル配布
\item IEEE MTT-S Student Research Grants ：機材支援プログラム
\end{itemize}

\subsection{SMR を買う}

\subsubsection{評価キット}

\begin{itemize}
\item \textbf{Qorvo}: 5G n77/n79 バンドの SMR フィルタ．
\$1,000 前後．Bragg 構造は公表されていないが，機密契約 (NDA) を結べば
設計資料を得られる場合あり
\item \textbf{Infineon}: SMR 評価キット．\$1,500 から
\end{itemize}

\subsubsection{研究用カスタム}

Fraunhofer IAF（ドイツ），CEA-Leti（フランス），VTT（フィンランド）で
研究用 SMR ウェハをカスタム作製可能．一枚 \$5,000--\$20,000．
これは DN 境界設計を自分で指定できるので最も clean．

\subsubsection{国内リソース}

村田製作所，TDK の研究所は SMR 開発を行っている．大学経由で
共同研究依頼すればサンプル提供の可能性あり．\textbf{これは共同研究
必須だが，提案書を「多モード真空揺らぎの機械共振器での
キャラクタリゼーション」と書けば RF 研究者に響きやすい}．

\subsection{入手のしやすさランキング}

\begin{center}
\begin{tabular}{ll}
\toprule
方法 & 入手難度 \\
\midrule
Digi-Key/Mouser でパッケージ品 & ◎ 即日 \\
評価キット（オンライン購入）& ○ 1--4 週 \\
メーカー直 bare die 請求 & △ 1--3 ヶ月 \\
研究用カスタム SMR & △ 6 ヶ月 \\
国内大学・企業連携 & △ 6 ヶ月--1 年 \\
\bottomrule
\end{tabular}
\end{center}

\subsection{個人で始めるなら}

\textbf{最初の一歩}：Digi-Key で 2.4\,GHz WiFi FBAR フィルタ
（例: Qorvo QPQ1903，\$8）と 5G n79 SMR フィルタ（例: Qorvo QM77020，
\$25）を各 5 個買う．合計 ¥25,000．\textbf{今日発注，今週届く}．

これとnanoVNA V2 Plus4 があれば測定 1（モードスペクトル確認）は
即座に可能．

\section{費用と時間}

\subsection{段階別予算}

\begin{center}
\begin{tabular}{lll}
\toprule
段階 & 予算 & 期間 \\
\midrule
Stage 0: nanoVNA ＋商用品 & ¥30,000 & 1 週間 \\
Stage 1: 自宅で温度スイープ & ¥80,000 & 1 ヶ月 \\
Stage 2: 本格 VNA 借用 ＋ Lamb shift & ¥300,000 & 3 ヶ月 \\
Stage 3: 液体窒素冷却 & ¥500,000 & 6 ヶ月 \\
Stage 4: 研究用カスタム SMR & ¥2,000,000 & 1 年 \\
\bottomrule
\end{tabular}
\end{center}

\subsection{Stage 0 の詳細（今すぐ開始可能）}

\begin{center}
\begin{tabular}{ll}
\toprule
項目 & 価格 \\
\midrule
Qorvo QPQ1903 (FBAR, 2.4 GHz) $\times 5$ & ¥6,000 \\
Qorvo QM77020 (SMR, 3.9 GHz) $\times 5$ & ¥18,000 \\
SMA コネクタ評価基板 $\times 2$ & ¥4,000 \\
nanoVNA V2 Plus4 & ¥25,000 \\
SMA 校正キット & ¥8,000 \\
SMA ケーブル ¥3,000 $\times 3$ & ¥9,000 \\
半田ごて・消耗品 & ¥10,000 \\
\midrule
\textbf{合計} & \textbf{¥80,000} \\
\bottomrule
\end{tabular}
\end{center}

1 週間あれば全て揃い，翌週末にはモードスペクトル測定が可能．
\textbf{これは軸 D（音響アナログ）と並ぶ即時着手案}．

\subsection{Stage 1 の詳細}

Stage 0 に加えて：
\begin{itemize}
\item ペルチェ温度制御ステージ ¥30,000
\item 熱電対と温度計 ¥5,000
\item 断熱ボックス（発泡スチロール改造）¥2,000
\item 電源と制御 Arduino ¥8,000
\end{itemize}

合計 ¥125,000．温度スイープ差分法が可能になり，測定 4 を実施できる．

\section{軸 D との関係}

冒頭で触れた「軸 D が null でも軸 A は死なない」理論は本ガイドにも
適用される．SMR/FBAR ルートは：

\begin{itemize}
\item 軸 D（空気音響管）より遥かに clean な DN 実装
\item しかし完全な量子極限ではない（室温 Q $\sim$ 3000）
\item 液体窒素や希釈冷凍機で量子に近づけられる
\end{itemize}

\begin{tcolorbox}[colback=blue!5,colframe=blue!50!black,title=推奨実行順序]
\begin{enumerate}
\item \textbf{今週}：軸 D 音響管（¥50,000）と SMR/FBAR Stage 0 (¥80,000)
を\textbf{同時発注}．総額 ¥130,000
\item \textbf{来週}：両方組み立て．モード構造確認
\item \textbf{1 ヶ月}：どちらかで偶数倍音欠落が統計的に見えれば第 1
成果．arXiv プレプリント
\item \textbf{3 ヶ月}：SMR/FBAR Stage 2 で温度スイープ Lamb shift
差を探索
\item \textbf{6 ヶ月}：結果次第で BAW（水晶）ルート，または共同研究で
cQED ルートに進む
\end{enumerate}
\end{tcolorbox}

\section{プラトー予想の検証：厚さシリーズ実験}

ここまでの測定（1--6）は「DN と NN でモード構造由来の Lamb shift
差が存在するか」を問う．\textbf{これは符号の検証}である．しかし
Wright Brothers プロジェクトの最大の賭けどころは符号ではなく，
\textbf{真空エネルギー密度の距離依存性が $1/d^4$ から
定数プラトーに遷移するか}である．これが正しければ，DN 斥力を
マクロスケールにスケールアップする道が開ける．肉眼・ガレージで
浮揚が見える可能性はここに集約されている．

\subsection{プラトー予想とは何か}

標準 QED 予測：真空エネルギー密度 $\rho_{\rm vac}(d) \propto 1/d^4$．
距離を 10 倍にすると効果は 10000 分の 1 になる．これが $1/a^4$ の壁．

Wright Brothers 仮説：\textbf{ある臨界距離 $a_c$ を超えると
$\rho_{\rm vac}$ が定数化}する．原因は位相的保護（$K_1(\mathbb Z[1/2])$
の winding number による），または Euler 積切除が
「体積効果」に転化するため．

\begin{center}
\begin{tabular}{ll}
\toprule
モデル & $\rho_{\rm vac}(d)$ \\
\midrule
標準 QED & $C \cdot d^{-4}$ \\
中間的新物理 & $C \cdot d^{-n}$, $n < 4$ \\
Wright Brothers プラトー & $C \cdot d^{-4}$ for $d < a_c$; const for $d > a_c$ \\
\bottomrule
\end{tabular}
\end{center}

$a_c$ の大きさは理論的に決まっていない．nm オーダー（プラトー発見
しても実用価値小）から mm オーダー（革命的）まで桁違いに予測が
開く．\textbf{だから実験するしかない}．

\subsection{厚さシリーズ測定の原理}

異なる厚さ $d$ を持つ SMR を複数用意し，各々の Lamb shift
差 $\Delta f(d) = f^{\rm NN}(d) - f^{\rm DN}(d)$ を測定する．

標準予測（$\rho \propto d^{-4}$）の下では：
\begin{equation}
  \Delta f(d) \propto d \cdot \rho \cdot (\text{結合}) \propto d^{-3}
\end{equation}
（より正確な導出はプロジェクト論文 \texttt{topological\_casimir.tex}
を参照．$\Delta f$ は $d$ の負冪則で減衰する）

プラトー予想の下では：
\begin{equation}
  \Delta f(d) \propto \begin{cases}
    d^{-3} & d < a_c \\
    d^{+1} & d > a_c
  \end{cases}
\end{equation}
$d > a_c$ では\textbf{厚さを増やすほど効果が増大}．通常の物理では
有り得ない振る舞い．

\begin{tcolorbox}[colback=yellow!5,colframe=orange!60!black,title=判定基準]
$\log\Delta f$ vs $\log d$ プロットにおいて，
\begin{itemize}
\item 全領域で傾き $-3$ → 標準 QED 支持，プラトー無し
\item 傾きが $-3$ より浅い（$-1$ とか $0$）→ 新物理，ただしプラトー
の手前
\item 傾きが\textbf{正}に反転 → プラトー発見，反転点が $a_c$
\end{itemize}
傾き $-3$ からの有意な逸脱を $3\sigma$ 以上で検出することが目標．
\end{tcolorbox}

\subsection{必要サンプル範囲}

\textbf{これが最大の難問}．商用 SMR は厚さが製造時に固定される．

商用品で入手可能な「事実上の厚さシリーズ」：

\begin{center}
\begin{tabular}{lll}
\toprule
周波数 & 対応厚さ (AlN) & 入手性 \\
\midrule
900\,MHz (LTE B8) & $\sim 5.8\,\mu$m & 商用品 \\
1.8\,GHz (LTE B3) & $\sim 2.9\,\mu$m & 商用品 \\
2.4\,GHz (WiFi) & $\sim 2.2\,\mu$m & 商用品 \\
3.9\,GHz (5G n77) & $\sim 1.3\,\mu$m & 商用品 \\
5.5\,GHz (WiFi 6) & $\sim 0.95\,\mu$m & 商用品 \\
28\,GHz (5G mmWave) & $\sim 0.19\,\mu$m & 商用品（高価）\\
60\,GHz (WiGig) & $\sim 0.09\,\mu$m & 研究用 \\
\bottomrule
\end{tabular}
\end{center}

\textbf{$5.8\,\mu$m から $0.19\,\mu$m まで 30 倍のレンジ}を商用品で
カバーできる．これでも $1/d^3$ スケーリングの検証は可能
（最小と最大で 27000 倍の Lamb shift 差が期待される）．

商用レンジを超える必要がある場合：

\begin{itemize}
\item \textbf{より薄い方向}: 研究用 mmWave SMR（60--200 GHz）
  \begin{itemize}
  \item Fraunhofer IAF, CEA-Leti, imec で製作依頼可能
  \item 1 枚 \$5,000--\$20,000
  \item 厚さ 50--200 nm
  \end{itemize}
\item \textbf{より厚い方向}: カスタム低周波 SMR
  \begin{itemize}
  \item 厚さ 10--100\,$\mu$m の AlN 膜は標準プロセスを超える
  \item LiNbO$_3$ 単結晶板（\$500/枚）で代替可能，厚さ 100\,$\mu$m--1\,mm
  \item ただし LiNbO$_3$ と AlN では音速が違うので正規化が必要
  \end{itemize}
\item \textbf{水晶 BAW との統合}: \texttt{guide\_baw\_quartz\_ja.tex}
で扱った水晶 BAW（$d = 100\,\mu$m--1\,mm）を同じシリーズに組み込む
\end{itemize}

\textbf{理想的な厚さレンジ}: 100\,nm から 1\,mm （$10^4$ 倍）．
この範囲でプラトー遷移が起きれば必ず見える．

\subsection{プロトコル}

\paragraph{Phase A: 商用品での予備測定（¥200,000, 2 ヶ月）}

\begin{enumerate}
\item 上表の商用 SMR/FBAR ペアを 6 周波数帯（900 MHz--5.5 GHz）で
全て揃える
\item 各ペアについて測定 1--5 を実施
\item $\log\Delta f$ vs $\log d$ をプロット
\item 傾きをフィット．$-3$ からの逸脱を統計的に評価
\item 逸脱があれば Phase B へ．無ければ標準 QED 支持として論文化
\end{enumerate}

\paragraph{Phase B: レンジ拡張（¥1,000,000, 6 ヶ月）}

\begin{enumerate}
\item 28 GHz mmWave SMR を追加（最薄側の拡張）
\item LiNbO$_3$ 基板で低周波側を拡張
\item 水晶 BAW（$d\sim$ mm）を統合測定系に組み込む
\item Phase A のプロットを 2 桁拡張
\item プラトー点 $a_c$ の上下限を絞る
\end{enumerate}

\paragraph{Phase C: カスタムウェハで決定測定（¥5,000,000, 1 年）}

\begin{enumerate}
\item Phase A/B で疑わしい $a_c$ 近傍に測定点を集中
\item Fraunhofer IAF に $d$ を 10 ポイント等間隔（対数スケール）で
作成依頼
\item 同じウェハ・同じ Bragg 設計で $d$ だけを変えた「厚さだけが違う」
サンプル群を作る
\item 系統誤差の大部分がキャンセルされる
\item プラトー発見の決定証拠を取得
\end{enumerate}

\subsection{系統誤差との戦い}

厚さシリーズ測定の最大の敵は「厚さ以外のパラメータが変わる」こと．

\textbf{厚さと共に変わるもの}：
\begin{itemize}
\item 基本周波数 $f_0 = v/(2d)$ or $v/(4d)$
\item Bragg 反射鏡の設計中心（$d$ に応じて調整される）
\item 電極の厚さ（通常 $d/10$ 程度で揃える）
\item Q 値（周波数依存性あり）
\item 測定系の校正（周波数帯で異なる）
\end{itemize}

\textbf{対策}：
\begin{enumerate}
\item 全ての量を $f_0$ で正規化し，次元解析的にスケール則を導出
\item 同じメーカー・同じプロセスのサンプルで統一
\item 測定系の校正は周波数ごとに厳密に実施
\item スケール不変な組合せ $\Delta f / f_0$ をプロット
\item Phase C の「厚さだけ変える」カスタムでクロスチェック
\end{enumerate}

$\Delta f/f_0$ は，標準 QED では $f_0$ に依存しない定数（自己相似的
$1/d^4$ のため）．プラトー領域では $f_0$ に\textbf{増加方向の依存性}を
持つ．これが clean な signature．

\subsection{プラトー発見の場合の次の一手}

もし Phase A--C のどこかでプラトー発見に至ったら（期待されるが
保証されない）：

\begin{enumerate}
\item \textbf{$a_c$ の精密決定}：プラトー遷移の距離を $\pm 10\%$ で
決める
\item \textbf{プラトー密度 $\rho_\infty$ の絶対測定}：単位体積あたりの
真空エネルギー差
\item \textbf{材料依存性の検証}：異なる圧電材料（AlN, LiNbO$_3$,
水晶, ZnO）でプラトー特性が共通か
\item \textbf{マクロプロトタイプ設計}：$a_c$ とプラトー密度が
分かれば，肉眼レベルの装置を設計できる
\item \textbf{理論側の展開}：$a_c$ と他の物理定数（$\hbar$, $c$,
プランク長，Casimir 定数）の関係を調べる．何かと一致すれば
新物理の基礎
\end{enumerate}

\begin{tcolorbox}[colback=green!5,colframe=green!60!black,title=プラトー発見が意味するもの]
\begin{itemize}
\item ダークエネルギー問題の解決候補（$\rho_\infty$ が観測値と
一致すれば）
\item マクロカシミール浮揚技術の誕生
\item QED の mode-sum 描像の根本的修正
\item Wright Brothers プロジェクトの理論的中核の実験的証明
\item ノーベル賞級の結果（誇張ではなく）
\end{itemize}
これが全て ¥200k の Phase A から始まる．
\end{tcolorbox}

\subsection{Phase A の具体部品リスト}

既存の Stage 0 (¥80,000) に加えて：

\begin{center}
\begin{tabular}{ll}
\toprule
追加項目 & 価格 \\
\midrule
Qorvo QPQ1294 (LTE B8, 900 MHz FBAR) $\times 3$ & ¥6,000 \\
Qorvo QM11007 (LTE B3, 1.8 GHz SMR) $\times 3$ & ¥9,000 \\
Broadcom ACFF-1024 (2.4 GHz FBAR) $\times 3$ & ¥9,000 \\
Qorvo QM77020 (3.9 GHz SMR, Stage 0 から) & (0) \\
Qorvo QM78017 (5.5 GHz SMR) $\times 3$ & ¥15,000 \\
Broadcom AFEM-9032 (5.2 GHz FBAR) $\times 3$ & ¥12,000 \\
広帯域 VNA レンタル（1 ヶ月 Keysight N5230C）& ¥80,000 \\
校正キット拡張（低周波用）& ¥20,000 \\
ケーブル・アダプタ追加 & ¥10,000 \\
\midrule
\textbf{追加合計} & \textbf{¥161,000} \\
\textbf{Stage 0 と合わせて総計} & \textbf{¥241,000} \\
\bottomrule
\end{tabular}
\end{center}

※ 具体的な型番はデータシート確認必要．表は概算．

\subsection{Phase A で期待される結果とその意味}

\begin{center}
\begin{tabular}{lll}
\toprule
観測される傾き & 解釈 & 次ステップ \\
\midrule
$-3 \pm 0.1$ & 標準 QED 確認 & 論文化．プラトーは nm 以下 or 存在せず \\
$-3 \pm 0.3$ & あいまい & Phase B で精度向上 \\
$-2$ 程度 & 何らかの新物理 & Phase B 必須 \\
$-1$ 以下 & プラトー接近 & Phase B 緊急 \\
$0$ or $+$ & \textbf{プラトー発見} & 緊急論文化，Phase C \\
\bottomrule
\end{tabular}
\end{center}

\subsection{正直な期待値}

\textbf{プラトー発見の事前確率は低い}．標準 QED は極めて精密に
検証されており，$1/d^4$ からの有意な逸脱は多くの既存実験に矛盾を
引き起こす可能性がある．一方で，\textbf{DN 境界での大規模 $d$
スイープは誰もやっていない}．既存の Casimir 測定は DD のみ，
nm スケールのみ．

したがってプラトー発見は「低確率・高リターン」の賭け．ただし：
\begin{itemize}
\item プラトー外れでも「傾き $-3$ の高精度確認」で論文 1 本
\item 符号反転（Section 8, 測定 1--6）の結果は別途確実に得られる
\item 厚さシリーズ自体が新規測定で，価値がある
\end{itemize}

\textbf{期待値 = 確実な成果 + 小確率での革命}．これが Phase A を
実行する合理的な理由．

\section{論文化戦略}

\subsection{第 1 論文: モード構造の直接観測}

タイトル案: ``Direct observation of Dirichlet--Neumann mode-parity
asymmetry in commercial bulk acoustic wave resonators''

内容:
\begin{itemize}
\item 商用 FBAR (NN) と SMR (DN) の $S_{11}$ 測定
\item 偶数倍音の欠落を統計的有意性込みで提示
\item Bragg 反射鏡帯域内外での ON/OFF
\item 歴史的文脈（Boyer 1974, Euler 積切除）に触れる
\end{itemize}

投稿先: \textit{Appl.\ Phys.\ Lett.} または \textit{J.\ Appl.\ Phys.}

新規性：産業界では既知の事実を物理学視点で初めて定量化．

\subsection{第 2 論文: Lamb shift 差の測定}

タイトル案: ``Boundary-mode-parity dependent vacuum Lamb shift
in bulk acoustic resonators''

内容:
\begin{itemize}
\item 温度スイープ差分法で FBAR/SMR の Lamb shift 差抽出
\item Boyer (1974) 予言 $\zeta_{\neg 2}(-3) = -7/120$ との定量比較
\item QED の mode-sum 予測との整合性 or 逸脱
\end{itemize}

投稿先: \textit{Phys.\ Rev.\ Lett.}

新規性：DN カシミール機械版の初の定量検証．

\subsection{第 3 論文: 理論的再解釈}

タイトル案: ``Commercial RF filters as realizations of Boyer's
1974 vacuum-repulsion prediction: a number-theoretic perspective''

Wright Brothers プロジェクトの哲学全体を示す．Euler 積切除，
Spec(Z)，位相的保護予想，そして SMR/FBAR 実験との接続．

投稿先: \textit{Found.\ Phys.} または arXiv のみ

\section{落とし穴と対策}

\subsection{Bragg 帯域の限界}

前述の通り Bragg 反射鏡は狭帯域．対策：
\begin{itemize}
\item メーカーから Bragg 設計データを入手（NDA 必要な場合あり）
\item 複数の中心周波数を持つ SMR 品を比較
\item 帯域内外の比較自体を検証手段にする
\end{itemize}

\subsection{パッケージング損失}

商用 SMR は金属缶パッケージ．缶自体が機械振動とカップルして
Q を劣化させる．\textbf{Bare die}を入手するのが理想だが，
パッケージ品でも基礎測定は可能．

\subsection{電極の非対称性}

上下電極の厚さ差は DN vs NN の純粋比較を崩す．実際の SMR は電極が
厚く，追加的な質量負荷で下面が「ほぼクランプ」になる．これは
我々の目的には好都合（DN を強める方向）．

\subsection{横モード（スプリアス）}

FBAR/SMR は縦モード以外に横波・Lamb 波なども持つ．これらは
メインの縦モードと結合してフィーチャーを作る．厳密な mode
identification が必須．メーカーのアプリケーションノートに詳細．

\subsection{熱的モード占有数}

室温 300\,K, $f = 2.6$\,GHz では $k_B T/h f \approx 2400$．モードは
強く熱占有されており，完全な「真空」状態ではない．これが Lamb
shift 測定で問題になるかは理論的に詰める必要あり．低温化で
解消．

\subsection{再現性}

同じロット内のバラつきは $\pm 0.1\%$ 程度．複数サンプルの統計平均で
抑える．ロットが異なると数 \% ずれるので，同一ロットから両タイプを
入手するのが理想（メーカー交渉必要）．

\section{Wright Brothers プロジェクト内での位置付け}

\subsection{マスタープランへの接続}

\texttt{wright\_brothers\_master\_plan.tex} の Stage 2
（3D Casimir Repulsion）の\textbf{最も安価で早い実装経路}が
SMR/FBAR 比較．これが成功すれば：

\begin{enumerate}
\item Boyer 予言の初の実験的証拠
\item $\zeta_{\neg 2}(-3)$ 予測の定量検証
\item Stage 3（Arakelov 幾何）への実験的動機付け
\item 研究資金調達における実績（他ルートへの拡張予算獲得）
\end{enumerate}

\subsection{他ルートとの統合}

\begin{itemize}
\item \textbf{軸 D 音響}: 古典領域の基礎検証．SMR と並行
\item \textbf{軸 A BAW (水晶)}: 次段階．Stage 2 が終わったら Stage 3
\item \textbf{軸 A cQED}: 最終検証．OtaNano 等での共同研究
\item \textbf{軸 C OPO}: 長期理論プログラム
\end{itemize}

\subsection{失敗しても得られるもの}

もし SMR/FBAR で偶数倍音欠落が見えなかった場合：
\begin{itemize}
\item Bragg 反射鏡の実効性の再評価（技術的貢献）
\item 産業品の物理特性データの学術公開（RF 工学に貢献）
\item 次ルートに進むための除外データ
\end{itemize}

ネガティブ結果でも論文化可能．

\section{まとめ}

SMR/FBAR ルートは Wright Brothers プロジェクトの DN カシミール
検証において現時点で最も現実的な経路である：

\begin{itemize}
\item ¥80,000 で Stage 0 開始可能
\item 商用品が構造上 DN 境界を実装済み
\item 大学の RF/MEMS 研究室と連携しやすい
\item 産業データの物理学的再解釈という明確な新規性
\item Boyer (1974) が 50 年前に予言した効果が，既に我々のポケットの
中にある可能性
\end{itemize}

軸 D（音響アナログ）と同時発注し，両者を並行して進めるのが
最適戦略．どちらかで先に成果が出れば，他方の位置付けが明確になる．

\vspace{1em}
\hrule
\vspace{0.5em}
\small
\noindent
関連文書：
\begin{itemize}
\item \texttt{guide\_garage\_scaling\_ja.tex}（4 軸の全体俯瞰）
\item \texttt{guide\_baw\_quartz\_ja.tex}（水晶 BAW ルート）
\item \texttt{cqed\_lambda4\_proposal.tex}（cQED ルート）
\item \texttt{boyer\_reinterpretation\_ja.tex}（理論背景）
\item \texttt{wright\_brothers\_master\_plan.tex}（全体計画）
\end{itemize}

\vspace{0.5em}
\noindent
主要参考文献：
\begin{itemize}
\item T.~H.~Boyer, Phys.\ Rev.\ A \textbf{9}, 2078 (1974)
\item R.~Aigner \textit{et al.}, Proc.\ IEEE Ultrason.\ Symp.\ (2003)
\item R.~Ruby \textit{et al.}, IEEE Microw.\ Mag.\ \textbf{16}, 60 (2015)
\item K.~Hashimoto (ed.), \textit{RF Bulk Acoustic Wave Filters for
Communications} (Artech House, 2009)
\item Y.~Chu \textit{et al.}, Science \textbf{358}, 199 (2017)
\end{itemize}

\end{document}
